
KV cache memory grows linearly with context length in transformer-based language models, consuming substantial memory during inference. Compression methods including H2O's heavy-hitter eviction and StreamingLLM's attention sink retention achieve significant memory reduction with reported minimal accuracy degradation on aggregate benchmarks. Quantization provides an orthogonal compression axis, reducing precision from 16-bit to 8-bit or 4-bit representations.

These methods share a fundamental assumption: one compression strategy fits all tasks. H2O applies identical heavy-hitter eviction ratios regardless of whether the task requires multi-document reasoning or single-document extraction. StreamingLLM uses fixed window sizes whether processing sequential narrative or precise retrieval tasks. This uniformity may sacrifice accuracy that task-appropriate strategy selection could recover.

The literature lacks systematic quantification of whether tasks genuinely respond differently to compression. Do some tasks tolerate aggressive eviction while others collapse? Does quantization affect summarization differently than code completion? Without empirical answers, practitioners resort to trial-and-error or conservative choices.

This work addresses this gap through rigorous quantification of task-dependent KV cache compression behavior. The contributions are:

\begin{enumerate}
\item \textbf{Task-compression clustering.} Gap statistic analysis demonstrates that 21 LongBench tasks form 3 distinct clusters (k* = 3) in their response to 6 compression configurations, validated with B = 500 bootstrap samples. The clusters align with interpretable task categories: Multi-doc QA + Code, Single-doc QA + Few-shot, and Summarization + Synthetic.
\end{enumerate}

\begin{enumerate}
\item \textbf{Attention entropy as task discriminator.} First-100-token attention entropy varies significantly across 6 task domains (F = 38.05, p = 2.92 $\times 10^{-26}, \eta^2$ = 0.522), indicating that 52\% of entropy variance is explained by task domain.
\end{enumerate}

\begin{enumerate}
\item \textbf{Identification of evaluation limitations.} The hypothesis that high-entropy tasks tolerate eviction better could not be validated (p = 0.326) due to baseline accuracy of 6.7\% rendering retention measurement unreliable. This identifies a methodological gap for future work.
\end{enumerate}

\begin{enumerate}
\item \textbf{Empirical foundation for task-conditioned compression.} These findings establish that task-aware strategy selection is grounded in measurable structure, motivating development of task-conditioned compression systems.
\end{enumerate}
